\documentclass[11pt,oneside]{article}

\pdfminorversion=7
\usepackage{QuizStyle}

\hypersetup{
    pdftitle={20261026 Quiz for MATH210 Applied Complex Variables},
    pdfauthor={Jungwoo Kim},
    pdfsubject={Quiz for MATH210 Applied Complex Variables},
    pdfcreator={LaTeX}
}

\begin{document}

\quiztitle{20261026 Quiz}

\begin{problem}{1}
Find all $z\in\C$ such that $\sin z=1$.
\end{problem}

\begin{solution}
Using the exponential definition of the complex sine, set $w=e^{iz}\ne0$. Then
\[
    \sin z=\frac{e^{iz}-e^{-iz}}{2i}=1
    \quad\Longleftrightarrow\quad
    w^2-2iw-1=(w-i)^2=0.
\]
Thus $e^{iz}=i$. Writing $z=x+iy$, we have
\[
    e^{iz}=e^{-y}(\cos x+i\sin x)=i.
\]
Comparison of moduli gives $e^{-y}=1$, hence $y=0$; comparison of arguments then gives $x=\pi/2+2\pi k$ for $k\in\Z$. Therefore all solutions are
\[
    \boxed{z=\frac{\pi}{2}+2\pi k,\qquad k\in\Z.}
\]
\end{solution}

\begin{problem}{2}
Find $a, b\in\R$ such that
\[
    \frac{(\sqrt3+i)^{20}}{(1-i)^{10}}=a+bi.
\]
\end{problem}

\begin{solution}
The polar forms of the two bases are
\[
    \sqrt3+i=2e^{i\pi/6},
    \qquad
    1-i=\sqrt2\,e^{-i\pi/4}.
\]
Consequently, by De Moivre's theorem,
\[
\begin{aligned}
    \frac{(\sqrt3+i)^{20}}{(1-i)^{10}}
    &=\frac{2^{20}e^{i10\pi/3}}{2^5e^{-i5\pi/2}}
      =2^{15}e^{i35\pi/6}\\
    &=2^{15}\left(\frac{\sqrt3}{2}-\frac{i}{2}\right)
      =2^{14}(\sqrt3-i).
\end{aligned}
\]
Hence
\[
    \boxed{a=2^{14}\sqrt3,\qquad b=-2^{14}.}
\]
\end{solution}

\begin{problem}{3}
Find an entire function $f(z)$ such that
\[
    \RePart f(z)=x^2-y^2+y,
    \qquad z=x+iy,
\]
and $f(0)=i$.
\end{problem}

\begin{solution}
Since
\[
    \RePart(z^2)=x^2-y^2,
    \qquad
    \RePart(-iz)=y,
\]
the polynomial $z^2-iz$ has the prescribed real part. An imaginary constant does not change its real part, and $f(0)=i$ fixes that constant. Thus
\[
    \boxed{f(z)=z^2-iz+i.}
\]
This function is entire, has the required real part, and satisfies $f(0)=i$.
\end{solution}

\begin{problem}{4}
Evaluate
\[
    \lim_{R\to\infty}\int_{C_R}
    \frac{2z^2-1}{z^4+z^2+1}\,\dd z,
\]
where $C_R$ is the circle $|z-1|=R$ with counterclockwise orientation.
\end{problem}

\begin{solution}
For $z\in C_R$, the reverse triangle inequality gives
\[
    R-1\le|z|\le R+1.
\]
For sufficiently large $R$, in particular for $R\ge5$, the denominator satisfies
\[
\begin{aligned}
    |z^4+z^2+1|
    &\ge |z|^4-|z|^2-1\\
    &\ge (R-1)^4-(R+1)^2-1>0,
\end{aligned}
\]
while $|2z^2-1|\le 2(R+1)^2+1$. Since the length of $C_R$ is $2\pi R$, the ML inequality yields
\[
\begin{aligned}
    \left|\int_{C_R}\frac{2z^2-1}{z^4+z^2+1}\,\dd z\right|
    &\le\frac{2\pi R(2(R+1)^2+1)}{(R-1)^4-(R+1)^2-1}
    \longrightarrow0.
\end{aligned}
\]
Therefore
\[
    \boxed{\lim_{R\to\infty}\int_{C_R}
    \frac{2z^2-1}{z^4+z^2+1}\,\dd z=0.}
\]
\end{solution}

\begin{problem}{5}
\begin{problemsubparts}
    \item Evaluate
    \[
        \int_C\frac{e^{z^2}}{z}\,\dd z,
    \]
    where $C$ is the circle $|z|=1$, oriented counterclockwise.
    \item Use part (a) to evaluate
    \[
        \int_0^{2\pi}e^{\cos 2\theta}\cos(\sin 2\theta)\,\dd\theta.
    \]
\end{problemsubparts}
\end{problem}

\begin{solution}
\begin{problemsubparts}
    \item The function $g(z)=e^{z^2}$ is entire, and $C$ is positively oriented with $0$ in its interior. By the Cauchy integral formula,
    \[
        \boxed{\int_C\frac{e^{z^2}}{z}\,\dd z
        =2\pi i\,g(0)=2\pi i.}
    \]
    \item Parametrize $C$ by $z=e^{i\theta}$, $0\le\theta\le2\pi$, so that $\dd z/z=i\,\dd\theta$. Part (a) becomes
    \[
        2\pi i=i\int_0^{2\pi}e^{e^{2i\theta}}\,\dd\theta.
    \]
    Euler's formula gives the essential decomposition
    \[
        e^{e^{2i\theta}}
        =e^{\cos 2\theta}
          \bigl(\cos(\sin 2\theta)+i\sin(\sin 2\theta)\bigr).
    \]
    Dividing by $i$ and taking real parts, we obtain
    \[
        \boxed{\int_0^{2\pi}e^{\cos 2\theta}
          \cos(\sin 2\theta)\,\dd\theta=2\pi.}
    \]
\end{problemsubparts}
\end{solution}

\begin{problem}{6}
Evaluate
\[
    \int_C\frac{e^z+\cos z}{(z-i)^{2025}}\,\dd z,
\]
where $C$ is the boundary of the square with vertices
$4\exp(i2k\pi/4)$ ($k=0, 1, 2, 3$), oriented clockwise.
\end{problem}

\begin{solution}
The vertices are $4, 4i, -4, -4i$, so $i$ lies inside $C$. Set $g(z)=e^z+\cos z$, which is entire. The exponent $2025$ in the denominator corresponds to the $2024$th derivative in the generalized Cauchy integral formula. Since $2024$ is divisible by $4$,
\[
    g^{(2024)}(z)=e^z+\cos z,
    \qquad
    g^{(2024)}(i)=e^i+\cos i.
\]
The given clockwise orientation reverses the usual sign. Therefore
\[
    \boxed{\int_C\frac{e^z+\cos z}{(z-i)^{2025}}\,\dd z
    =-\frac{2\pi i}{2024!}(e^i+\cos i).}
\]
\end{solution}

\begin{problem}{7}
Let $P(z)$ be a nonconstant polynomial. Describe the image of the set
\[
    A:=\{z\in\C\mid P(z)\ne0\}
\]
under the map $f(z)=\exp(1/P(z))$.
\end{problem}

\begin{solution}
By the fundamental theorem of algebra, every $w\in\C$ is attained by $P$: indeed, the nonconstant polynomial $P(z)-w$ has a root. Restricting to $A$ therefore gives
\[
    P(A)=\C\setminus\{0\},
    \qquad
    \left\{\frac1{P(z)}\mid z\in A\right\}=\C\setminus\{0\}.
\]
The exponential never vanishes, so $f(A)\subseteq\C\setminus\{0\}$. Conversely, take any $w\ne0$. A logarithm $L$ of $w$ can be chosen nonzero: if $w\ne1$, every logarithm is nonzero; if $w=1$, choose $L=2\pi i$. Since $L\ne0$, there exists $z\in A$ with $1/P(z)=L$, and then $f(z)=e^L=w$. Hence
\[
    \boxed{f(A)=\C\setminus\{0\}.}
\]
\end{solution}

\begin{problem}{8}
Consider the power series
\[
    \sum_{n=1}^{\infty}a_nz^n,
    \qquad
    a_n=\begin{cases}
        n & \text{if $n$ is even},\\
        1/n & \text{if $n$ is odd}.
    \end{cases}
\]
\begin{problemsubparts}
    \item Find the radius of convergence of the series.
    \begin{hint}
    You may evaluate $\displaystyle\lim_{n\to\infty}a_n^{1/n}$.
    \end{hint}
    \item Find all $z$ at which the series converges.
\end{problemsubparts}
\end{problem}

\begin{solution}
\begin{problemsubparts}
    \item Along the even and odd indices, respectively,
    \[
        a_n^{1/n}=n^{1/n}\longrightarrow1,
        \qquad
        a_n^{1/n}=n^{-1/n}\longrightarrow1.
    \]
    Thus $\lim_{n\to\infty}a_n^{1/n}=1$. The Cauchy--Hadamard formula gives
    \[
        \boxed{R=\frac1{\limsup_{n\to\infty}|a_n|^{1/n}}=1.}
    \]
    \item By the root test, the series converges absolutely for $|z|<1$ and diverges for $|z|>1$. If $|z|=1$, its even-index terms satisfy
    \[
        |a_{2k}z^{2k}|=2k\longrightarrow\infty,
    \]
    so the terms do not tend to zero; the series diverges at every boundary point. Therefore the exact convergence set is
    \[
        \boxed{\{z\in\C\mid |z|<1\}.}
    \]
\end{problemsubparts}
\end{solution}

\begin{problem}{9}
Show that the series
\[
    \sum_{n=1}^{\infty}\frac{e^{inz}}{n^2}
\]
converges uniformly on $\ImPart z\ge0$.
\end{problem}

\begin{solution}
For $z=x+iy$ with $y=\ImPart z\ge0$, we have
\[
    \left|\frac{e^{inz}}{n^2}\right|
    =\frac{|e^{inx-ny}|}{n^2}
    =\frac{e^{-ny}}{n^2}\le\frac1{n^2}.
\]
Since $\sum_{n=1}^{\infty}1/n^2$ converges and its bound is independent of $z$, the Weierstrass M-test proves that
\[
    \boxed{\sum_{n=1}^{\infty}\frac{e^{inz}}{n^2}
      \text{ converges uniformly on }\{z\in\C\mid\ImPart z\ge0\}.}
\]
\end{solution}

\Needspace{15\baselineskip}
\begin{problem}{10}
Determine whether each statement is true or false, and justify your answer.
\begin{problemsubparts}
    \item There exists an analytic function $f(z)$ with
    $\RePart f(z)=x^3-y^3$, where $z=x+iy$.
    \item If $f(z)$ is entire and $|f(z)|\ge C>0$ for all $z\in\C$,
    then $f$ is a constant function.
\end{problemsubparts}
\end{problem}

\begin{solution}
\begin{problemsubparts}
    \item \textbf{False.} The real part $u$ of an analytic function is harmonic. But for $u(x, y)=x^3-y^3$,
    \[
        u_{xx}+u_{yy}=6x-6y,
    \]
    which is not identically zero on any nonempty open set. Thus $u$ cannot be the real part of an analytic function on an open domain.

    \item \textbf{True.} Since $|f(z)|\ge C>0$, the entire function $f$ never vanishes. Hence $h(z)=1/f(z)$ is entire and
    \[
        |h(z)|=\frac1{|f(z)|}\le\frac1C
        \qquad(z\in\C).
    \]
    Liouville's theorem makes $h$ constant. Its constant value is nonzero, so $f=1/h$ is constant as well.
\end{problemsubparts}
\end{solution}

\end{document}
